\documentclass{article}
\usepackage{graphicx}
\begin{document}
\title{Tiny Transformers for Forecasting}
\maketitle

\begin{abstract}
We ask whether a small transformer matches larger baselines on forecasting (RQ1) and how sensitive it is to its training settings (RQ2).
\end{abstract}

\section{Introduction}\label{sec:intro}
Attention models are common in forecasting \cite{vaswani2017attention}.
RQ1: does a small model match larger ones? RQ2: how sensitive is it to hyperparameters?

\section{Related work}\label{sec:related}
Attention was introduced by \cite{vaswani2017attention} and refined by \cite{devlin2019bert}.
A recent report \cite{smith2021ghost} disagrees; see also \cite{nobody2020}.

\section{Method}\label{sec:method}
The architecture is shown in Figure~\ref{fig:missing}. We minimise the loss in Eq.~\ref{eq:loss}.
\begin{equation}\label{eq:loss}
\mathcal{L} = \sum_t (y_t - \hat y_t)^2
\end{equation}
\begin{figure}
% generated-by: manual
\includegraphics{figures/arch.pdf}
\caption{Model architecture.}\label{fig:arch}
\end{figure}

\section{Experimental setup}\label{sec:setup}
We train with a learning rate of $3\times10^{-4}$ and batch size 64, using the Adam optimizer.
Details are in Appendix~\ref{app:hparams}.

\section{Results}\label{sec:results}
Table~\ref{tab:main} reports the main results; our model reaches 92.4 F1 on the benchmark.
Figure~\ref{fig:curve} shows training curves and Figure~\ref{fig:qual} shows examples.
\begin{table}
% source: code:results/table2.csv
\caption{Main results.}\label{tab:main}
\begin{tabular}{lr}
Baseline & 0.871 \\
Ours & 0.902 \\
\end{tabular}
\end{table}
\begin{table}
\caption{Ablation.}\label{tab:ablation}
\begin{tabular}{lr}
No attention & 0.842 \\
\end{tabular}
\end{table}
\begin{figure}
% generated-by: code:plot_curve.py
\includegraphics{figures/curve.png}
\caption{Training curves.}\label{fig:curve}
\end{figure}
\begin{figure}
\includegraphics{figures/qual.png}
\caption{Qualitative examples.}\label{fig:qual}
\end{figure}

\section{Conclusion}\label{sec:concl}
A small model is competitive (RQ1) and moderately sensitive to settings (RQ2).

\appendix
\section{Extra results}\label{app:extra}
Additional curves are omitted for space.

\section{Hyperparameters}\label{app:hparams}
\begin{tabular}{ll}
Learning rate & $1\times10^{-3}$ \\
Batch size & 128 \\
Optimizer & Adam \\
\end{tabular}

\bibliography{refs}
\end{document}
